\subsection{Свойства функций, имеющих пределы: единственность, предел модуля, предельные переходы в неравенствах, теорема о зажатой функции}

\textbf{Свойства функций, имеющих предел}\\
Всюду далее считаем, что пределы рассматриваются при $x \to x_0$ (или $x \to \infty$).

\textbf{Теорема.} Если функция $f(x)$ имеет предел, то он единственен.
\begin{tcolorbox}[title=1. Единственность предела]
    \textbf{Теорема.} Функция не может иметь двух различных пределов в одной точке.
    
    \textbf{Доказательство (от противного):}
    Предположим, что $\lim_{x\to a} f(x) = A$ и $\lim_{x\to a} f(x) = B$, причем $A \neq B$.
    
    1. Выберем $\varepsilon = \frac{|A - B|}{2} > 0$.
    2. По определению предела:
    \begin{itemize}
        \item Для $A$: $\exists \delta_1 > 0: \forall x \in \mathring{U}_{\delta_1}(a) \implies |f(x) - A| < \varepsilon$.
        \item Для $B$: $\exists \delta_2 > 0: \forall x \in \mathring{U}_{\delta_2}(a) \implies |f(x) - B| < \varepsilon$.
    \end{itemize}
    3. Возьмем $\delta = \min(\delta_1, \delta_2)$. Для любого $x \in \mathring{U}_\delta(a)$ выполняются оба неравенства.
    4. Оценим расстояние между $A$ и $B$ через неравенство треугольника:
    \[ |A - B| = |(A - f(x)) + (f(x) - B)| \leq |A - f(x)| + |f(x) - B| \]
    Подставим оценки:
    \[ |A - B| < \varepsilon + \varepsilon = 2\varepsilon = 2 \cdot \frac{|A - B|}{2} = |A - B| \]
    5. Получили $|A - B| < |A - B|$, что невозможно.
    Противоречие доказывает, что $A = B$.
\end{tcolorbox}

\textbf{Теорема.} Если $\lim f(x) = A$, то $\lim |f(x)| = |A|$.
\begin{tcolorbox}[title=2. Предел модуля]
	
	\textbf{Доказательство:}
	Используем неравенство треугольника вида $||a| - |b|| \leq |a - b|$.
	\begin{itemize}
		\item Зададим $\forall \varepsilon > 0$.
		\item Так как $f(x) \to A$, то $\exists \delta: \forall x \in \mathring{U}_\delta(x_0) \Rightarrow |f(x) - A| < \varepsilon$.
		\item Тогда:
		$$ ||f(x)| - |A|| \leq |f(x) - A| < \varepsilon $$
		\item Это и означает, что $|f(x)| \to |A|$.
	\end{itemize}
	\textit{Важно: Обратное утверждение верно только в случае $A=0$.}
\end{tcolorbox}

\textbf{Теорема.} Пусть $\lim f(x) = A$ и $\lim g(x) = B$.
Если в некоторой проколотой окрестности точки $x_0$ выполняется $f(x) \leq g(x)$, то $A \leq B$.
	
\begin{tcolorbox}[title=3. Предельный переход в неравенствах]
	\textbf{Доказательство (от противного):}
	\begin{enumerate}
		\item Допустим, что $A > B$.
		\item Возьмем $\varepsilon = \frac{A - B}{2} > 0$.
		\item Для этого $\varepsilon$ найдутся $\delta_1$ и $\delta_2$. Возьмем $\delta = \min(\delta_1, \delta_2)$. В общей окрестности $\mathring{U}_\delta(x_0)$:
		$$ |f(x) - A| < \varepsilon \Rightarrow f(x) > A - \varepsilon = A - \frac{A-B}{2} = \frac{A+B}{2} $$
		$$ |g(x) - B| < \varepsilon \Rightarrow g(x) < B + \varepsilon = B + \frac{A-B}{2} = \frac{A+B}{2} $$
		\item Получаем цепочку: $g(x) < \frac{A+B}{2} < f(x)$, то есть $g(x) < f(x)$.
		\item Это противоречит условию $f(x) \leq g(x)$. Значит, предположение неверно, и $A \leq B$.
	\end{enumerate}
	\textit{Замечание: Строгое неравенство $f(x) < g(x)$ при переходе к пределу может стать нестрогим (например, $1/x^2 > 0$, но предел равен 0).}
\end{tcolorbox}

\textbf{Теорема.} Пусть в некоторой окрестности точки $x_0$ выполняется:
$$ g(x) \leq f(x) \leq h(x) $$
И пусть $\lim_{x\to x_0} g(x) = \lim_{x\to x_0} h(x) = A$.
Тогда $\lim_{x\to x_0} f(x) = A$.
\begin{tcolorbox}[title=4. Теорема о зажатой функции (о двух милиционерах)]	
	\textbf{Доказательство:}
	\begin{enumerate}
		\item Запишем определение предела для крайних функций. $\forall \varepsilon > 0 \exists \delta$:
		\begin{itemize}
			\item $A - \varepsilon < g(x) < A + \varepsilon$
			\item $A - \varepsilon < h(x) < A + \varepsilon$
		\end{itemize}
		\item Используем двойное неравенство из условия (для $x$ из общей окрестности):
		$$ A - \varepsilon < g(x) \leq f(x) \leq h(x) < A + \varepsilon $$
		\item Уберем промежуточные члены:
		$$ A - \varepsilon < f(x) < A + \varepsilon $$
		$$ |f(x) - A| < \varepsilon $$
		\item Это и означает по определению, что $\lim f(x) = A$.
	\end{enumerate}
\end{tcolorbox}
